\introChapter{การนำเสนอผลลัพธ์และการทำแผนที่ (Cartography and Data Visualization)}

\begin{description}[leftmargin=2.5cm, style=nextline]
    \item[รหัสวิชา] 4121701
    \item[ชื่อวิชา] คณิตศาสตร์สำหรับคอมพิวเตอร์
    \item[หน่วยกิต] 3 (3-0-6)
    \item[บทที่] 7
    \item[ชื่อบท] การนำเสนอผลลัพธ์และการทำแผนที่ (Cartography and Data Visualization)
    \item[จำนวนชั่วโมง] 6 ชั่วโมง
\end{description}

\vspace{1em}
\noindent\textbf{วัตถุประสงค์การเรียนรู้ประจำบท}

\noindent เมื่อนักศึกษาเรียนบทนี้จบแล้ว นักศึกษาจะสามารถ:

\begin{enumerate}
    \item อธิบายหลักการออกแบบแผนที่และการจัดลำดับความสำคัญของข้อมูลภาพ (Visual Hierarchy) ได้
    \item เลือกใช้สัญลักษณ์ (Symbology) ที่เหมาะสมกับข้อมูลเชิงพื้นที่ ทั้งในด้านสี รูปร่าง และขนาด ได้
    \item จำแนกประเภทข้อมูลและเลือกวิธีการจำแนกข้อมูล (Classification Methods) ที่เหมาะสม ได้
    \item ออกแบบและจัดทำเลย์เอาต์แผนที่ (Map Layout) ที่สื่อสารข้อมูลเชิงพื้นที่ได้อย่างมีประสิทธิภาพ
    \item สื่อสารผลวิเคราะห์ข้อมูลเชิงพื้นที่ให้อ่านง่ายและถูกต้องตามหลักสากล
\end{enumerate}

\vspace{1em}
\noindent\textbf{หัวข้อที่สอน}

\begin{enumerate}
    \item หลักการออกแบบแผนที่ (Visual Hierarchy)
    \item การเลือกสัญลักษณ์ (Symbology): สี รูปร่าง ขนาด
    \item การจำแนกข้อมูล (Classification Methods): Natural Breaks, Quantile, Equal Interval
    \item การจัดทำเลย์เอาต์แผนที่ (Map Layout)
    \item การสื่อสารข้อมูลเชิงพื้นที่ให้อ่านง่ายและถูกต้องตามหลักสากล
\end{enumerate}

\vspace{1em}
\noindent\textbf{สื่อการสอน}
\begin{itemize}
    \item สไลด์ประกอบการสอน
    \item ซอฟต์แวร์ GIS สำหรับสร้างแผนที่
    \item แบบฝึกหัดท้ายบท
\end{itemize}
